and therefore, for every @@M0@@:
@@M1@@

By (7.1.1), the @@M2@@ are isobaric polynomials of weight @@M3@@ in the @@M4@@, with @@M5@@. For example, let us give the values of @@M6@@ and @@M7@@:
@@M8@@

\sgaStableTarget{sga6.e05.para.7-2}7.2. Let @@M9@@ be a ring of characteristic @@M10@@; suppose given additive maps @@M11@@ from @@M12@@ to @@M13@@, such that @@M14@@. We then define a pre-@@M15@@-structure on @@M16@@ by putting:
\sgaStableTarget{sga6.e05.eq.7-2-1}@@M17@@
Multiplicativity follows from the additivity of the @@M18@@. There is therefore an equivalence between giving a pre-@@M19@@-structure on @@M20@@ and giving the corresponding operations @@M21@@.

\sgaStableTarget{sga6.e05.para.7-3}7.3. It is clear from formula (7.1.1) that the @@M22@@ commute with @@M23@@-homomorphisms. If @@M24@@ are pre-@@M25@@-rings of characteristic zero, it is necessary and sufficient that a homomorphism @@M26@@ commute with the @@M27@@ for it to be a @@M28@@-homomorphism: this follows from 7.2, the @@M29@@ being expressed as polynomials in the @@M30@@.

\sgaStableTarget{sga6.e05.ex.7-4}\textbf{Examples 7.4.} i) For @@M31@@, endowed with its @@M32@@-ring structure, we have:
@@M33@@
Consequently, the @@M34@@ are all equal to the identity map of @@M35@@.

ii) Let us work in the category @@M36@@ of @@M37@@-algebras (1.8). The
